% GENERATED FILE. Edit canonical lessons, then run npm run generate:books.
% canonical-source-hash: fnv1a64:8e4a1cfe45d44bb9
% canonical-lessons: HI-C449-siyar, HI-C449-jugnu, HI-C449-jhigur, HI-C449-carna, HI-C449-carana

\chapter{Night in the Fields}
\label{ch:hi-night-in-the-fields}
\hlchaptermodality{449}

\begin{chapteropening}
\textbf{By the end of this chapter:} \emph{I can describe night in the fields: jackals, fireflies and crickets, and cattle out grazing.}

Complete the last lesson of chapter 449: I can describe night in the fields: jackals, fireflies and crickets, and cattle out grazing.
\end{chapteropening}

\section[\texorpdfstring{\hi{सियार} (siyār)}{सियार (siyār)}]{\hi{सियार} (siyār) — a jackal}

\label{lesson:HI-C449-siyar}



\begin{quote}
Before the new one: say the Hindi for to peck, then the Hindi for to prick.
\end{quote}

\subsection*{You'll want to know: \hi{सियार}}
\textbf{\hi{सियार}} — \emph{siyār} — \textquotedblleft{}a jackal\textquotedblright{}.

\begin{quote}
\hi{रात} \hi{को} \hi{खेतों} \hi{में} \hi{सियार} \hi{घूमते} \hi{हैं।}
\end{quote}

\begin{quote}
\emph{rāt ko khetoṁ meṁ siyār ghūmte haiṁ.}
\end{quote}

\begin{quote}
Jackals roam the fields at night.
\end{quote}

Masculine.

\begin{cousinweb}
\textbf{\hi{सियार}} (\emph{siyār}) is a tadbhava of Sanskrit \emph{śṛgāla}, \textquotedblleft{}a jackal\textquotedblright{} (CDIAL 12578): the g between vowels has dropped and the vowels have run together.
\end{cousinweb}

\begin{grammarlens}[title={What you've built}]
One more word for getting around.
\end{grammarlens}

\subsection*{Your turn}
\begin{itemize}
\raggedright
  \item \emph{Say it:} \emph{siyār}
  \item \emph{Say it:} \emph{siyār}, once more
  \item \emph{Say it:} \emph{rāt ko khetoṁ meṁ siyār ghūmte haiṁ.} — Jackals roam the fields at night.
  \item \emph{Recall:} say the Hindi for to peck, then the Hindi for to prick, then say \emph{siyār} again
\end{itemize}

\begin{tcolorbox}[breakable,colback=teal!4,colframe=teal!35!black,title={Before you move on}]
What does \hi{सियार} mean? (\textquotedblleft{}A jackal\textquotedblright{}.) Say it once more.
\end{tcolorbox}
\section[\texorpdfstring{\hi{जुगनू} (jugnū)}{जुगनू (jugnū)}]{\hi{जुगनू} (jugnū) — a firefly}

\label{lesson:HI-C449-jugnu}



\begin{quote}
Before the new one: say the Hindi for to prick, then the Hindi for a jackal.
\end{quote}

\subsection*{You'll want to know: \hi{जुगनू}}
\textbf{\hi{जुगनू}} — \emph{jugnū} — \textquotedblleft{}a firefly\textquotedblright{}.

\begin{quote}
\hi{रात} \hi{को} \hi{खेतों} \hi{में} \hi{जुगनू} \hi{दिखते} \hi{हैं।}
\end{quote}

\begin{quote}
\emph{rāt ko khetoṁ meṁ jugnū dikhte haiṁ.}
\end{quote}

\begin{quote}
You can see fireflies in the fields at night.
\end{quote}

Masculine.

\begin{cousinweb}
\textbf{\hi{जुगनू}} (\emph{jugnū}) is a tadbhava of Sanskrit \emph{jyotiriṅgaṇa}, \textquotedblleft{}a firefly\textquotedblright{}, literally a moving light (CDIAL 5298), which lists Hindi \emph{jugnū}.
\end{cousinweb}

\begin{grammarlens}[title={What you've built}]
One more word for getting around.
\end{grammarlens}

\subsection*{Your turn}
\begin{itemize}
\raggedright
  \item \emph{Say it:} \emph{jugnū}
  \item \emph{Say it:} \emph{jugnū}, once more
  \item \emph{Say it:} \emph{rāt ko khetoṁ meṁ jugnū dikhte haiṁ.} — You can see fireflies in the fields at night.
  \item \emph{Recall:} say the Hindi for to prick, then the Hindi for a jackal, then say \emph{jugnū} again
\end{itemize}

\begin{tcolorbox}[breakable,colback=teal!4,colframe=teal!35!black,title={Before you move on}]
What does \hi{जुगनू} mean? (\textquotedblleft{}A firefly\textquotedblright{}.) Say it once more.
\end{tcolorbox}
\section[\texorpdfstring{\hi{झींगुर} (jhī̃gur)}{झींगुर (jhī̃gur)}]{\hi{झींगुर} (jhī̃gur) — a cricket (insect)}

\label{lesson:HI-C449-jhigur}



\begin{quote}
Before the new one: say the Hindi for a jackal, then the Hindi for a firefly.
\end{quote}

\subsection*{You'll want to know: \hi{झींगुर}}
\textbf{\hi{झींगुर}} — \emph{jhī̃gur} — \textquotedblleft{}a cricket (insect)\textquotedblright{}.

\begin{quote}
\hi{रात} \hi{को} \hi{झींगुर} \hi{बहुत} \hi{शोर} \hi{करते} \hi{हैं।}
\end{quote}

\begin{quote}
\emph{rāt ko jhī̃gur bahut śor karte haiṁ.}
\end{quote}

\begin{quote}
Crickets make a lot of noise at night.
\end{quote}

Masculine.

\begin{cousinweb}
\textbf{\hi{झींगुर}} (\emph{jhī̃gur}) goes back to \emph{*jhiṅgura}, a form CDIAL reconstructs (CDIAL 5387) and glosses as a cockroach; it lists Hindi \emph{jhĩgur}, \textquotedblleft{}a cricket\textquotedblright{}.
\end{cousinweb}

\begin{grammarlens}[title={What you've built}]
One more word for getting around.
\end{grammarlens}

\subsection*{Your turn}
\begin{itemize}
\raggedright
  \item \emph{Say it:} \emph{jhī̃gur}
  \item \emph{Say it:} \emph{jhī̃gur}, once more
  \item \emph{Say it:} \emph{rāt ko jhī̃gur bahut śor karte haiṁ.} — Crickets make a lot of noise at night.
  \item \emph{Recall:} say the Hindi for a jackal, then the Hindi for a firefly, then say \emph{jhī̃gur} again
\end{itemize}

\begin{tcolorbox}[breakable,colback=teal!4,colframe=teal!35!black,title={Before you move on}]
What does \hi{झींगुर} mean? (\textquotedblleft{}A cricket (insect)\textquotedblright{}.) Say it once more.
\end{tcolorbox}
\section[\texorpdfstring{\hi{चरना} (carnā)}{चरना (carnā)}]{\hi{चरना} (carnā) — to graze}

\label{lesson:HI-C449-carna}



\begin{quote}
Before the new one: say the Hindi for a firefly, then the Hindi for a cricket.
\end{quote}

\subsection*{You'll want to know: \hi{चरना}}
\textbf{\hi{चरना}} — \emph{carnā} — \textquotedblleft{}to graze\textquotedblright{}.

\begin{quote}
\hi{मैदान} \hi{में} \hi{गायें} \hi{चर} \hi{रही} \hi{थीं।}
\end{quote}

\begin{quote}
\emph{maidān meṁ gāyẽ car rahī thīṁ.}
\end{quote}

\begin{quote}
Cows were grazing in the field.
\end{quote}

Past (it …): \hi{वह} \hi{भैंस} \hi{दिन} \hi{भर} \hi{चरी}, \emph{vah bhaĩs din bhar carī}.

\begin{cousinweb}
\textbf{\hi{चरना}} (\emph{carnā}) is a tadbhava of Sanskrit \emph{carati}, \textquotedblleft{}moves\textquotedblright{} and also \textquotedblleft{}grazes\textquotedblright{} (CDIAL 4686), which lists Hindi \emph{carnā}, \textquotedblleft{}to move about, graze\textquotedblright{}.
\end{cousinweb}

\begin{grammarlens}[title={What you've built}]
One more word for getting around.
\end{grammarlens}

\subsection*{Your turn}
\begin{itemize}
\raggedright
  \item \emph{Say it:} \emph{carnā}
  \item \emph{Say it:} \emph{carnā}, once more
  \item \emph{Say it:} \emph{maidān meṁ gāyẽ car rahī thīṁ.} — Cows were grazing in the field.
  \item \emph{Recall:} say the Hindi for a firefly, then the Hindi for a cricket, then say \emph{carnā} again
\end{itemize}

\begin{tcolorbox}[breakable,colback=teal!4,colframe=teal!35!black,title={Before you move on}]
What does \hi{चरना} mean? (\textquotedblleft{}To graze\textquotedblright{}.) Say it once more.
\end{tcolorbox}
\section[\texorpdfstring{\hi{चराना} (carānā)}{चराना (carānā)}]{\hi{चराना} (carānā) — to take to graze; to graze (cattle)}

\label{lesson:HI-C449-carana}



\begin{quote}
Before the new one: say the Hindi for a cricket, then the Hindi for to graze.
\end{quote}

\subsection*{You'll want to know: \hi{चराना}}
\textbf{\hi{चराना}} — \emph{carānā} — \textquotedblleft{}to take to graze; to graze (cattle)\textquotedblright{}.

\begin{quote}
\hi{एक} \hi{लड़का} \hi{नदी} \hi{के} \hi{किनारे} \hi{बकरियाँ} \hi{चरा} \hi{रहा} \hi{था।}
\end{quote}

\begin{quote}
\emph{ek laṛkā nadī ke kināre bakriyā̃ carā rahā thā.}
\end{quote}

\begin{quote}
A boy was grazing goats on the river bank.
\end{quote}

Past, with \hi{ने} (he or she …): \hi{उसने} \hi{भेड़ें} \hi{चराईं}, \emph{usne bheṛẽ carāīṁ}.

\begin{cousinweb}
\textbf{\hi{चराना}} (\emph{carānā}) is the causative partner of \hi{चरना}, from Sanskrit \emph{carati}, \textquotedblleft{}moves, grazes\textquotedblright{} (CDIAL 4686), which lists Hindi \emph{carānā}, \textquotedblleft{}to graze\textquotedblright{} a herd.
\end{cousinweb}

\begin{grammarlens}[title={What you've built}]
One more word for getting around.
\end{grammarlens}

\subsection*{Your turn}
\begin{itemize}
\raggedright
  \item \emph{Say it:} \emph{carānā}
  \item \emph{Say it:} \emph{carānā}, once more
  \item \emph{Say it:} \emph{ek laṛkā nadī ke kināre bakriyā̃ carā rahā thā.} — A boy was grazing goats on the river bank.
  \item \emph{Recall:} say the Hindi for a cricket, then the Hindi for to graze, then say \emph{carānā} again
\end{itemize}

\begin{tcolorbox}[breakable,colback=teal!4,colframe=teal!35!black,title={Before you move on}]
What does \hi{चराना} mean? (\textquotedblleft{}To take to graze; to graze (cattle)\textquotedblright{}.) Say it once more.
\end{tcolorbox}
